\subsection{Planteamiento del problema}

El uso de la tecnología para manejar los datos en los centros de salud es uno de los retos más grandes que existen hoy en día. Los sistemas de información médicos han cambiado mucho, pasando de las carpetas de papel a programas avanzados que buscan conectar las consultas de los doctores, las cuentas de cobro y los procesos administrativos en un solo lugar. Sin embargo, este cambio no pasa de la misma forma en todos los consultorios. En muchos centros de atención básica, los datos de los pacientes todavía están regados por todas partes, guardados en herramientas que no se comunican entre sí y dependiendo de que el personal copie todo a mano. Este desorden con la información no es un detalle simple, sino que es la causa principal de que el trabajo sea lento, se corran riesgos con las leyes de privacidad y el consultorio no pueda atender a más personas. Entender por qué pasa esto y cómo solucionarlo es el motor de este trabajo de titulación.

A nivel mundial, los investigadores coinciden en que tener la información dividida en archivos sueltos causa una carga de trabajo enorme para el personal de salud. Por ejemplo, en un análisis detallado sobre cómo trabaja el personal en los consultorios, se descubrió que los puntos donde la información se corta son muy comunes, lo que obliga a los empleados a pasar los datos manualmente de un programa a otro \cite{antonacci2021process}. De la misma manera, se ha señalado que la falta de un orden claro en las tareas empeora la desconexión entre lo que el doctor hace en la consulta y lo que se anota en administración, un problema que causa tareas repetidas, equivocaciones al registrar enfermedades y demoras que metodologías de diagramación de procesos intentan solucionar \cite{PUFAHL2022102013}. Además, los estudios que miden cuánto tiempo se pierde llenando papeles confirman que los profesionales de la salud pasan gran parte de su día en tareas de oficina en lugar de atender pacientes \cite{Murad2024DocumentationBurden}. Esto demuestra que trabajar con la información de forma manual es imposible de mantener cuando el consultorio empieza a recibir más gente.

En Ecuador y América Latina, la situación tiene problemas específicos que complican aún más las cosas. En una encuesta realizada a más de 300 odontólogos ecuatorianos, se descubrió que el miedo a que los datos de los pacientes no estén seguros es la mayor barrera para usar tecnología, seguido por la falta de tiempo para aprender a usar sistemas nuevos \cite{CherrezOjeda2020ICTDentists}. Por la parte legal, el país cuenta con la Ley Orgánica de Protección de Datos Personales desde el año 2021, y los análisis demuestran que tanto los centros públicos como privados tienen serios problemas para cumplir con las seguridades técnicas y legales que exige esta norma para cuidar los datos médicos, que son considerados muy sensibles \cite{martinez2022proteccion}. Esto representa un peligro real de multas o sanciones para los consultorios pequeños que guardan información en archivos locales sin contraseñas ni registros de quién entra a verlos. Por si fuera poco, las investigaciones locales revelan que existen fallas constantes al llenar los formularios oficiales de odontología del Ministerio de Salud Pública, sobre todo porque los doctores y secretarias no reciben la capacitación adecuada para usar los códigos de diagnóstico requeridos \cite{uce2016caracteristicas}. Este vacío no se arregla solo con cursos, sino con un software que ya tenga las reglas incluidas.

El centro odontológico Bellidental, que está ubicado en la Provincia de Tungurahua dentro del cantón Baños, es un reflejo exacto de toda esta situación. Con una atención de unos 75 pacientes a la semana y más de 1 000 expedientes acumulados, la clínica funciona con un modelo manual donde los datos clínicos, las cuentas y las citas están completamente desconectados. Al revisar el trabajo diario mediante entrevistas con el personal, se encontró que las historias clínicas se llenan en archivos de Excel que copian el formato del Ministerio de Salud \cite{uce2016caracteristicas}, las citas se anotan en Google Calendar, los pagos se registran en hojas de papel firmadas a mano, las recetas se escriben con pluma y los recordatorios se envían escribiendo uno por uno en WhatsApp. Toda esta separación obliga a que la secretaria y el doctor funcionen como un puente humano, arrastrando datos de un lado a otro. Las consecuencias se notan todos los días, ya que la secretaria debe interrumpir la consulta del odontólogo para preguntarle qué códigos poner en el Excel, buscar el historial de un paciente entre miles de filas de cálculo toma demasiado tiempo, los pacientes no entienden la letra de las recetas médicas, hay pérdidas de dinero por citas olvidadas y no hay forma de saber si los tratamientos cuadran con los ingresos de la caja. A esto se suma que tener información sensible en una computadora sin seguridad expone al centro a problemas legales con la ley de protección de datos \cite{martinez2022proteccion}. En resumen, el verdadero problema de Bellidental es que su flujo de información carece de estandarización e integración, lo que aumenta el tiempo que se pierde en tareas de oficina por cada paciente, vuelve la operación diaria muy pesada y frena la capacidad de la clínica para atender a más personas y seguir creciendo.